\section{代码风格与静态检查}

用 \code{pylint} 4.1.3 做静态检查，标准取 PEP 8（行宽 79、4 空格缩进），配置文件 \code{pylintrc} 放在代码根目录。只关掉三类与作业场景无关的检查：\code{duplicate-code}（模板化的 dataclass 声明必然相似）、\code{too-few-public-methods}（实体本来就只有数据）、\code{too-many-lines}。参数上限放宽到 6 个、属性上限 14 个，因为 \code{Item} 有 13 个字段，这是领域模型本身的样子，拆开反而难读。

第一版可运行代码（不含测试）跑出 \textbf{9.85/10}，共 27 条告警，分布见\tabref{tab:pylintinit}。没有一条是逻辑错误类（\code{E} 级）告警，全是格式与写法问题——"能跑"和"好读"确实是两件事。

\begin{table}[htbp]
\centering
\caption{初稿 27 条告警的分布}\label{tab:pylintinit}
\zihao{-5}
\renewcommand{\arraystretch}{1.15}
\begin{tabularx}{\textwidth}{@{}L{5.0cm}L{1.2cm}L{3.6cm}Y@{}}
\toprule
告警 & 条数 & 典型位置 & 处理方式 \\
\midrule
\code{line-too-long} C0301 & 22 & 6 个文件，最长 83 字符 & 按语义折行，不缩排 \\
\code{trailing-newlines} C0305 & 2 & \code{item.py}、\code{moderation.py} & 删除文件末尾多余空行 \\
\code{unidiomatic-typecheck} C0123 & 1 & \code{storage.py:60} & \code{type(None)} 改为 \code{types.NoneType} \\
\code{too-many-return-statements} R0911 & 1 & \code{storage.py:54} & 重写为单一出口 \\
\code{unused-import} W0611 & 1 & \code{auth.py:12} & 删除未使用的 \code{NotFound} \\
\bottomrule
\end{tabularx}
\end{table}

行宽占 22 条，其中 15 条只超 1--4 个字符，删几个空格也能过，但那样只是把问题藏起来。按语义折行才是对的：函数调用的参数对齐到左括号后，\code{if} 的条件用括号包起来再按布尔项分行。

\begin{lstlisting}
# 修复前（83 字符）
self._journal.record(actor.user_id, action, item.item_id, True, item.title)

# 修复后：按"调用参数对齐"折行
self._journal.record(
    actor.user_id, action, item.item_id, True, item.title)
\end{lstlisting}

把测试代码一起纳入检查后重新跑，得分 \textbf{10.00/10}：

\lstinputlisting{listings/pylint_after.tex}

测试目录有三处文件级 \code{disable}，都写了理由：pytest 夹具互相引用是官方写法（\code{redefined-outer-name}）、用例名本身就是文档（\code{missing-function-docstring}）、有的夹具只负责铺数据（\code{unused-argument}）。除此之外没有其它豁免，源码目录一条 \code{disable} 都没有。

这次也吃了晚检查的亏：第一次跑 pylint 是写完 3000 行之后，27 条里 22 条是行宽，属于写的时候就该注意的问题，攒到最后改就是一次几百行的机械劳动。
